\chapter{Hyperbolic Isometries}
\label{ch:vii34-hyperbolic-isometries}

VII/33 identified \(PSL(2,\mathbb R)\) with the orientation-preserving isometry
group of the hyperbolic plane.  We now classify individual elements by their
fixed points and dynamics.

\section{Fixed-point equation}

Let
\[
A=
\begin{pmatrix}
a&b\\c&d
\end{pmatrix}
\in SL(2,\mathbb R).
\]
A fixed point satisfies
\[
z=\frac{az+b}{cz+d},
\]
or
\[
cz^2+(d-a)z-b=0.
\]

\begin{proposition}[Discriminant]\label{prop:vii34-discriminant}
The discriminant of the fixed-point equation is
\[
(a+d)^2-4
=
(\operatorname{tr}A)^2-4.
\]
\end{proposition}

Thus the absolute trace controls the qualitative dynamics.

\section{Elliptic, parabolic, hyperbolic}

\begin{definition}[Classification]\label{def:vii34-classification}
A nonidentity element represented by \(A\in SL(2,\mathbb R)\) is:
\[
\begin{array}{ll}
\text{elliptic} & |\operatorname{tr}A|<2,\\
\text{parabolic} & |\operatorname{tr}A|=2,\\
\text{hyperbolic} & |\operatorname{tr}A|>2.
\end{array}
\]
Because \(A\) and \(-A\) represent the same element, only the absolute trace is
well defined on \(PSL(2,\mathbb R)\).
\end{definition}

\section{Elliptic elements}

\begin{theorem}[Elliptic fixed point]\label{thm:vii34-elliptic}
An elliptic element has one fixed point in \(\mathbb H\) and no fixed point on
\(\partial_\infty\mathbb H\).  It is conjugate to a rotation fixing \(i\).
\end{theorem}

\begin{example}[Rotation about \(i\)]\label{ex:vii34-rotation}
Matrices
\[
R_\theta=
\begin{pmatrix}
\cos\theta&\sin\theta\\
-\sin\theta&\cos\theta
\end{pmatrix}
\]
fix \(i\).  Their action is elliptic unless they represent the identity.
\end{example}

\section{Parabolic elements}

\begin{theorem}[Parabolic normal form]\label{thm:vii34-parabolic}
A parabolic element has exactly one fixed point on the ideal boundary and is
conjugate to
\[
z\longmapsto z+1.
\]
\end{theorem}

The unique boundary fixed point is often called a parabolic fixed point.

\section{Hyperbolic elements}

\begin{theorem}[Hyperbolic normal form]\label{thm:vii34-hyperbolic}
A hyperbolic element has exactly two fixed points on the ideal boundary and is
conjugate to
\[
z\longmapsto e^\ell z
\]
for a unique number \(\ell>0\).
\end{theorem}

\begin{definition}[Axis]\label{def:vii34-axis}
The unique geodesic joining the two boundary fixed points of a hyperbolic element
is its axis.
\end{definition}

\begin{proposition}[Translation along the axis]\label{prop:vii34-axis-translation}
A hyperbolic element preserves its axis and translates every point on it by the
same hyperbolic distance \(\ell\).
\end{proposition}

\section{Translation length}

\begin{definition}[Translation length]\label{def:vii34-translation-length}
For an isometry \(g\), define
\[
\tau(g)
=
\inf_{z\in\mathbb H}d(z,gz).
\]
\end{definition}

\begin{theorem}[Three cases]\label{thm:vii34-tau-cases}
For nonidentity orientation-preserving isometries:
\[
\begin{array}{lll}
\text{elliptic:} & \tau(g)=0, & \text{attained in }\mathbb H,\\
\text{parabolic:} & \tau(g)=0, & \text{not attained in }\mathbb H,\\
\text{hyperbolic:} & \tau(g)>0, & \text{attained precisely on the axis}.
\end{array}
\]
\end{theorem}

\section{Trace and translation length}

\begin{theorem}[Trace-length relation]\label{thm:vii34-trace-length}
If \(g\) is hyperbolic and represented by \(A\), then
\[
|\operatorname{tr}A|
=
2\cosh\frac{\tau(g)}2.
\]
Equivalently,
\[
\tau(g)
=
2\operatorname{arcosh}\frac{|\operatorname{tr}A|}{2}.
\]
\end{theorem}

\begin{proof}
Conjugate to the normal form
\[
z\mapsto e^\ell z.
\]
A determinant-one representative has eigenvalues \(e^{\ell/2}\) and
\(e^{-\ell/2}\), so its trace has absolute value
\[
e^{\ell/2}+e^{-\ell/2}
=
2\cosh(\ell/2).
\]
The axis is the imaginary line and the displacement there is \(\ell\).
\end{proof}

\section{Boundary dynamics}

Let \(g\) be hyperbolic with boundary fixed points
\[
g^-,g^+.
\]

\begin{proposition}[North--south dynamics]\label{prop:vii34-north-south}
For every boundary point \(x\ne g^-\),
\[
g^n x\to g^+
\qquad(n\to+\infty),
\]
and for every \(x\ne g^+\),
\[
g^{-n}x\to g^-.
\]
\end{proposition}

The point \(g^+\) is attracting and \(g^-\) repelling.

\section{Parabolic dynamics}

A parabolic element has one boundary fixed point \(p\).  After conjugating
\(p\) to \(\infty\), the map is a translation \(z\mapsto z+b\).  Iterates move
toward the same ideal point in both forward and backward directions, though from
opposite sides in boundary coordinates.

\section{Elliptic order}

\begin{proposition}[Finite versus infinite order]\label{prop:vii34-elliptic-order}
An elliptic element conjugate to a rotation through angle \(\theta\) has finite
order exactly when
\[
\theta/\pi
\]
is rational.
\end{proposition}

\section{Conjugacy invariants}

\begin{proposition}[Type is conjugacy invariant]\label{prop:vii34-conjugacy}
Elliptic, parabolic, and hyperbolic type are invariant under conjugation in
\(PSL(2,\mathbb R)\).
\end{proposition}

\begin{proof}
Conjugation preserves fixed-point structure and preserves the trace of a matrix
representative up to the sign ambiguity already absorbed by absolute trace.
\end{proof}

\section{Orientation-reversing isometries}

The full isometry group also contains reflections in hyperbolic geodesics and
glide-reflection analogues.

\begin{remark}
The classification above concerns orientation-preserving isometries.  The
orientation-reversing side can be obtained by composing these elements with a
fixed reflection.
\end{remark}

\section{Why this matters for groups}

When many hyperbolic isometries are assembled into a discrete subgroup, their
fixed points, axes, translation lengths, and interactions control the geometry of
the quotient surface.  This is the entry point to Fuchsian groups.

\section{Exercises}
\begin{exercise}\label{exr:vii34-01}
Write the fixed-point equation for \(A(z)=(az+b)/(cz+d)\).
\end{exercise}
\begin{hint}
Cross multiply.
\end{hint}
\begin{solution}
\(cz^2+(d-a)z-b=0\).
\end{solution}

\begin{exercise}\label{exr:vii34-02}
What is its discriminant when \(\det A=1\)?
\end{exercise}
\begin{hint}
Simplify \((d-a)^2+4bc\).
\end{hint}
\begin{solution}
\((\operatorname{tr}A)^2-4\).
\end{solution}

\begin{exercise}\label{exr:vii34-03}
Classify an element with \(|\operatorname{tr}A|<2\).
\end{exercise}
\begin{hint}
Use the trace trichotomy.
\end{hint}
\begin{solution}
It is elliptic.
\end{solution}

\begin{exercise}\label{exr:vii34-04}
Classify \(|\operatorname{tr}A|=2\), nonidentity.
\end{exercise}
\begin{hint}
Use the trace trichotomy.
\end{hint}
\begin{solution}
It is parabolic.
\end{solution}

\begin{exercise}\label{exr:vii34-05}
Classify \(|\operatorname{tr}A|>2\).
\end{exercise}
\begin{hint}
Use the trace trichotomy.
\end{hint}
\begin{solution}
It is hyperbolic.
\end{solution}

\begin{exercise}\label{exr:vii34-06}
Why use absolute trace in \(PSL(2,\mathbb R)\)?
\end{exercise}
\begin{hint}
\(A\) and \(-A\) are identical projectively.
\end{hint}
\begin{solution}
The trace changes sign under \(A\mapsto-A\), while its absolute value does not.
\end{solution}

\begin{exercise}\label{exr:vii34-07}
Where is the fixed point of an elliptic isometry?
\end{exercise}
\begin{hint}
Look inside the half-plane.
\end{hint}
\begin{solution}
It has one fixed point in \(\mathbb H\).
\end{solution}

\begin{exercise}\label{exr:vii34-08}
How many ideal fixed points does a parabolic element have?
\end{exercise}
\begin{hint}
Discriminant zero.
\end{hint}
\begin{solution}
Exactly one.
\end{solution}

\begin{exercise}\label{exr:vii34-09}
How many ideal fixed points does a hyperbolic element have?
\end{exercise}
\begin{hint}
Discriminant positive.
\end{hint}
\begin{solution}
Exactly two.
\end{solution}

\begin{exercise}\label{exr:vii34-10}
Give a parabolic normal form.
\end{exercise}
\begin{hint}
Move its fixed point to infinity.
\end{hint}
\begin{solution}
\(z\mapsto z+1\).
\end{solution}

\begin{exercise}\label{exr:vii34-11}
Give a hyperbolic normal form.
\end{exercise}
\begin{hint}
Move its fixed points to \(0,\infty\).
\end{hint}
\begin{solution}
\(z\mapsto e^\ell z\), \(\ell>0\).
\end{solution}

\begin{exercise}\label{exr:vii34-12}
Define the axis of a hyperbolic isometry.
\end{exercise}
\begin{hint}
Join its two ideal fixed points.
\end{hint}
\begin{solution}
It is the unique geodesic with those two boundary endpoints.
\end{solution}

\begin{exercise}\label{exr:vii34-13}
What does a hyperbolic element do on its axis?
\end{exercise}
\begin{hint}
Use the normal form.
\end{hint}
\begin{solution}
It translates every point by the same positive distance.
\end{solution}

\begin{exercise}\label{exr:vii34-14}
Define translation length.
\end{exercise}
\begin{hint}
Take the infimum displacement.
\end{hint}
\begin{solution}
\(\tau(g)=\inf_z d(z,gz)\).
\end{solution}

\begin{exercise}\label{exr:vii34-15}
What is \(\tau(g)\) for elliptic \(g\)?
\end{exercise}
\begin{hint}
It has an interior fixed point.
\end{hint}
\begin{solution}
Zero, attained at the fixed point.
\end{solution}

\begin{exercise}\label{exr:vii34-16}
What is \(\tau(g)\) for parabolic \(g\)?
\end{exercise}
\begin{hint}
It has no interior fixed point.
\end{hint}
\begin{solution}
Zero, but the infimum is not attained in \(\mathbb H\).
\end{solution}

\begin{exercise}\label{exr:vii34-17}
What is \(\tau(g)\) for hyperbolic \(g\)?
\end{exercise}
\begin{hint}
Use its axis.
\end{hint}
\begin{solution}
It is positive and attained exactly on the axis.
\end{solution}

\begin{exercise}\label{exr:vii34-18}
State the trace--length relation.
\end{exercise}
\begin{hint}
Use \(\cosh(\tau/2)\).
\end{hint}
\begin{solution}
\(|\operatorname{tr}A|=2\cosh(\tau/2)\).
\end{solution}

\begin{exercise}\label{exr:vii34-19}
Which boundary fixed point of a hyperbolic element attracts forward iterates?
\end{exercise}
\begin{hint}
Use north--south dynamics.
\end{hint}
\begin{solution}
The attracting fixed point \(g^+\).
\end{solution}

\begin{exercise}\label{exr:vii34-20}
What happens under backward iterates?
\end{exercise}
\begin{hint}
Reverse the roles.
\end{hint}
\begin{solution}
They converge to the repelling fixed point \(g^-\) away from \(g^+\).
\end{solution}

\begin{exercise}\label{exr:vii34-21}
When does an elliptic rotation have finite order?
\end{exercise}
\begin{hint}
Require its angle to be rational relative to \(\pi\).
\end{hint}
\begin{solution}
Exactly when \(\theta/\pi\in\mathbb Q\).
\end{solution}

\begin{exercise}\label{exr:vii34-22}
Is isometry type invariant under conjugation?
\end{exercise}
\begin{hint}
Conjugation preserves fixed points/trace.
\end{hint}
\begin{solution}
Yes.
\end{solution}

\begin{exercise}\label{exr:vii34-23}
Give an orientation-reversing hyperbolic isometry type.
\end{exercise}
\begin{hint}
Think reflection.
\end{hint}
\begin{solution}
Reflection in a hyperbolic geodesic is the basic example.
\end{solution}

\begin{exercise}\label{exr:vii34-24}
State the bridge to VII/35.
\end{exercise}
\begin{hint}
Pass from single elements to discrete subgroups.
\end{hint}
\begin{solution}
VII/35 develops Fuchsian groups, fundamental domains, quotient surfaces, and the geometry encoded by discrete collections of these isometries.
\end{solution}

\section{Solved problem dossiers}
\begin{problem}[Discriminant problem]\label{prob:vii34-01}
Derive the trace discriminant of the fixed-point equation.
\end{problem}
\begin{solution}
The quadratic discriminant is \((d-a)^2+4bc=(a+d)^2-4(ad-bc)=(\operatorname{tr}A)^2-4\).
\end{solution}

\begin{problem}[Elliptic problem]\label{prob:vii34-02}
Explain why negative discriminant yields one interior fixed point.
\end{problem}
\begin{solution}
The two algebraic roots are complex conjugates; exactly one lies in the upper half-plane and the other in the lower half-plane.
\end{solution}

\begin{problem}[Parabolic problem]\label{prob:vii34-03}
Show \(z\mapsto z+1\) is parabolic.
\end{problem}
\begin{solution}
It is represented by \(\begin{smallmatrix}1&1\\0&1\end{smallmatrix}\), whose trace is \(2\), and its unique boundary fixed point is \(\infty\).
\end{solution}

\begin{problem}[Hyperbolic problem]\label{prob:vii34-04}
Show \(z\mapsto e^\ell z\) is hyperbolic.
\end{problem}
\begin{solution}
Represent it by \(\operatorname{diag}(e^{\ell/2},e^{-\ell/2})\); the trace is \(2\cosh(\ell/2)>2\), with fixed points \(0,\infty\).
\end{solution}

\begin{problem}[Axis problem]\label{prob:vii34-05}
Why is the imaginary axis the axis of \(z\mapsto e^\ell z\)?
\end{problem}
\begin{solution}
Its ideal endpoints are \(0\) and \(\infty\), and the unique geodesic with those endpoints is the imaginary axis.
\end{solution}

\begin{problem}[Translation-length problem]\label{prob:vii34-06}
Compute the displacement along the axis of \(z\mapsto e^\ell z\).
\end{problem}
\begin{solution}
Using vertical distance, \(d(iy,ie^\ell y)=|\log(e^\ell)|=\ell\).
\end{solution}

\begin{problem}[Trace-length problem]\label{prob:vii34-07}
Derive the trace--length formula.
\end{problem}
\begin{solution}
The determinant-one eigenvalues in normal form are \(e^{\ell/2},e^{-\ell/2}\); add them and identify \(\ell=\tau(g)\).
\end{solution}

\begin{problem}[Parabolic-infimum problem]\label{prob:vii34-08}
Why does a parabolic element have zero translation length without a minimizing point?
\end{problem}
\begin{solution}
For \(z\mapsto z+1\), points high in the half-plane have displacement tending to zero because the metric scale is \(1/y\), but no interior point is fixed.
\end{solution}

\begin{problem}[North-south problem]\label{prob:vii34-09}
Prove the boundary dynamics of a hyperbolic normal form.
\end{problem}
\begin{solution}
For \(z\mapsto e^\ell z\), every finite nonzero boundary point tends to \(\infty\) forward and to \(0\) backward; conjugation transports this behavior to arbitrary fixed points.
\end{solution}

\begin{problem}[Conjugacy problem]\label{prob:vii34-10}
Show each type is preserved by conjugacy.
\end{problem}
\begin{solution}
Conjugation transports fixed points bijectively and preserves the absolute trace of representatives, so the discriminant sign is unchanged.
\end{solution}

\begin{problem}[Elliptic-order problem]\label{prob:vii34-11}
Determine when a rotation about \(i\) has finite order.
\end{problem}
\begin{solution}
After conjugation to a disk rotation \(z\mapsto e^{i\phi}z\), finite order means \(e^{im\phi}=1\) for some integer \(m\), equivalently \(\phi/(2\pi)\in\mathbb Q\).
\end{solution}

\begin{problem}[Classification-synthesis problem]\label{prob:vii34-12}
Relate trace, fixed points, and translation length in all three cases.
\end{problem}
\begin{solution}
Elliptic: \(|\operatorname{tr}|<2\), one interior fixed point, zero attained length. Parabolic: \(|\operatorname{tr}|=2\), one ideal fixed point, zero unattained length. Hyperbolic: \(|\operatorname{tr}|>2\), two ideal fixed points, positive attained length.
\end{solution}

\section{Challenge problems}
\begin{problem}[Eigenvalue classification]\label{prob:vii34-ch01}
Rephrase the three types in terms of eigenvalues of a determinant-one representative.
\end{problem}
\begin{solution}
Elliptic eigenvalues are complex conjugates on the unit circle; parabolic has repeated eigenvalue \(\pm1\) with nontrivial Jordan form; hyperbolic has distinct real reciprocal eigenvalues \(\lambda,\lambda^{-1}\) with \(|\lambda|>1\).
\end{solution}

\begin{problem}[Axis uniqueness]\label{prob:vii34-ch02}
Why is the hyperbolic axis unique?
\end{problem}
\begin{solution}
The two distinct ideal fixed points determine a unique hyperbolic geodesic. Any invariant geodesic on which the element translates must have those endpoints, hence must be that geodesic.
\end{solution}

\begin{problem}[Displacement off axis]\label{prob:vii34-ch03}
Why should displacement of a hyperbolic element be strictly larger off its axis?
\end{problem}
\begin{solution}
Conjugate to a pure dilation. The axis realizes the orthogonal direction of translation; moving away introduces an additional transverse component. The hyperbolic distance formula shows the displacement increases with distance from the axis.
\end{solution}

\begin{problem}[Parabolic horocycles preview]\label{prob:vii34-ch04}
What Euclidean curves are naturally preserved by \(z\mapsto z+1\)?
\end{problem}
\begin{solution}
Horizontal lines \(y=\text{constant}\) are preserved. In hyperbolic geometry these are horocycles centered at the parabolic fixed point \(\infty\).
\end{solution}

\begin{problem}[Discrete-group bridge]\label{prob:vii34-ch05}
Why are elliptic finite order, parabolic cusps, and hyperbolic axes the three basic local behaviors in Fuchsian quotients?
\end{problem}
\begin{solution}
In a discrete subgroup, finite-order elliptics produce orbifold cone behavior, parabolics produce cusp ends, and hyperbolics identify along axes to create closed geodesic behavior. VII/35 organizes these phenomena group-theoretically.
\end{solution}

\section*{Bridge to VII/35}
VII/35 passes from individual hyperbolic isometries to discrete subgroups: Fuchsian groups, fundamental domains, quotient surfaces, cusps, and closed geodesic structures.
